\clearpage\pagestyle{empty}\thispagestyle{empty}
\addcontentsline{toc}{chapter}{Résumé}
\setstretch{1.05}\small

\noindent{\large\bfseries Résumé}\par\medskip

Ce rapport présente une étude comparative de pronostic de défaillance de roulements sur les données réelles du banc PRONOSTIA (FEMTO-ST, challenge IEEE PHM 2012). La fiabilité de la population, estimée par maximum de vraisemblance en tenant compte des roulements censurés, suit une loi de Weibull ($\beta$ = 1,85) : B10 = 125 min et B50 = 346 min, avec de larges intervalles de confiance. Un défaut est détecté par la règle des 3 $\sigma$ en médiane 7,8 min avant la panne, et un indicateur de santé est construit à partir d'énergies de bandes d'ondelettes. Sept méthodes de prédiction de la durée de vie résiduelle sont comparées : le processus de Wiener bayésien est le plus précis (RMSE de 16,4 min) et le LSTM le plus prudent. Des tests sur données simulées, des tests statistiques et une analyse d'explicabilité valident la démarche. Enfin, une maintenance prédictive fondée sur le Wiener réduit le coût par heure de 47 \% par rapport au correctif.

\textbf{Mots-clés :} fiabilité, loi de Weibull, vraisemblance, durée de vie résiduelle, processus de Wiener, apprentissage automatique, maintenance prédictive.

\par\medskip\noindent\rule{\linewidth}{0.4pt}\par\medskip
\begin{otherlanguage}{english}
\noindent{\large\bfseries Abstract}\par\medskip

This report presents a comparative study of bearing failure prognostics on the real PRONOSTIA dataset (FEMTO-ST, IEEE PHM 2012 challenge). The population reliability, estimated by maximum likelihood with right-censored bearings, follows a Weibull distribution ($\beta$ = 1.85): B10 = 125 min and B50 = 346 min, with wide confidence intervals. A first defect is detected by a 3-sigma rule a median of 7.8 min before failure, and a health indicator is built from wavelet band energies. Seven remaining useful life methods are compared: the Bayesian Wiener process is the most accurate (RMSE of 16.4 min) and the LSTM the most cautious. Tests on simulated data, statistical tests and an explainability analysis validate the approach. Finally, a predictive maintenance policy based on the Wiener model cuts the cost per hour by 47\% compared with corrective maintenance.

\textbf{Keywords:} reliability, Weibull distribution, likelihood, remaining useful life, Wiener process, machine learning, predictive maintenance.
\end{otherlanguage}

\par\medskip\noindent\rule{\linewidth}{0.4pt}\par\medskip
\resumearabe
